\chapter{Decision Problems}\label{chap:decisionproblems}

\firstwords{A fundamental aspect} of studying various models of computation is determining precisely what kinds of questions about that model can be answered algorithmically. For example, does there exist an algorithm that can tell us whether or not a deterministic finite automaton accepts a given input word? It seems obvious---after all, such an algorithm only needs to figure out whether the finite automaton, given the input word, would end up in an accepting state or not---but it's still something that we need to establish formally.

In the previous chapter, we introduced the notions of decidable and semidecidable languages. Using much of the same terminology we developed for reasoning about languages, we can talk more generally not just about languages but about \emph{problems} pertaining to our various models of computation, and we can formally prove which problems can be decided in much the same way as languages can be decided.

A \emph{decision problem} is a problem for which each input instance corresponds to either a ``yes" or a ``no" answer. At its core, a decision problem is a language, and the elements of the language are input instances that correspond to a ``yes" answer. Each element of the language is an encoding of whatever model of computation we're considering, in some cases with an input word given to that model: in the case of regular languages, we use encodings of finite automata, and so on.

If there exists a \emph{decision algorithm} that produces the correct ``yes" or ``no" answer for any instance of a given decision problem, and the algorithm halts on all inputs, then we say that the decision problem is \emph{decidable}. If no such decision algorithm exists, then we say that the decision problem is \emph{undecidable}.

\begin{remark}
Note that ``undecidable" doesn't mean the decision problem is impossible; ``undecidable" only means ``not decidable", as in ``no decision algorithm exists that \emph{always} halts and gives a yes/no answer".
\end{remark}

There are a number of common decision problems that we can ask about a model of computation $\textsf{X} \in \{\RE, \DFA, \NFA, \ENFA, \CFG, \PDA, \TM, \dots\}$, and these problems are summarized in Table~\ref{tab:commondecisionproblems}. In this chapter, we will get our first taste of ``programming" Turing machines to solve each of these decision problems. Note that we no longer need to construct a specific Turing machine for a particular decision problem, because (as we now know) the Church--Turing thesis tells us that any function that can be computed by an algorithm can also be computed by a Turing machine.

\begin{table}
\centering
\caption{Common decision problems and their definitions}
\label{tab:commondecisionproblems}
\begin{tabular}{l l}
\toprule
Membership	& $\mathit{A}_{\textsf{X}} = \{\langle \mathcal{B}, w \rangle \mid \mathcal{B} \text{ is an \textsf{X} that accepts input word } w\}$ \\
Emptiness		& $\mathit{E}_{\textsf{X}} = \{\langle \mathcal{B} \rangle \mid \mathcal{B} \text{ is an \textsf{X} and } L(\mathcal{B}) = \emptyset\}$ \\
Universality	& $\mathit{U}_{\textsf{X}} = \{\langle \mathcal{B} \rangle \mid \mathcal{B} \text{ is an \textsf{X} and } L(\mathcal{B}) = \Sigma^{*}\}$ \\
Equivalence	& $\mathit{EQ}_{\textsf{X}} = \{\langle \mathcal{B}, \mathcal{C} \rangle \mid \mathcal{B} \text{ and } \mathcal{C} \text{ are both \textsf{X} and } L(\mathcal{B}) = L(\mathcal{C})\}$ \\
Inclusion		& $\mathit{IN}_{\textsf{X}} = \{\langle \mathcal{B}, \mathcal{C} \rangle \mid \mathcal{B} \text{ and } \mathcal{C} \text{ are both \textsf{X} and } L(\mathcal{B}) \subseteq L(\mathcal{C})\}$ \\
\bottomrule
\end{tabular}
\end{table}

\begin{construction}
The study of computation doesn't end with decision problems. As Table~\ref{tab:typesofproblems} illustrates, a number of types of computational problems exist depending on what specifically we would like to know about a problem. Decision problems are the most common formulation, but I will eventually write a few paragraphs about these other problem types, as they will appear in later chapters.
\end{construction}

\begin{sidewaystable}
\centering
\caption{Types of computational problems}
\label{tab:typesofproblems}
\begin{tabular}{l >{\raggedright\arraybackslash}p{0.2\linewidth} l l l}
\toprule
\textbf{Type of problem}		& \textbf{Description}		& \textbf{Typical formulation}	& \textbf{Example} \\
\midrule
Decision problem			& Answer to every instance is either ``yes" or ``no" & \linebreakcell{l}{$f \from \Sigma^{*} \to \{0,1\}$ \\ $L = \{x \in \Sigma^{*} \mid f(x) = 1\}$} & \linebreakcell{l}{Primality testing: \\ $2 \rightarrow \text{``yes, prime"}$} \\
\addlinespace[0.5em]
Function problem			& Answer to every instance exists, but is not necessarily binary (i.e., ``yes" or ``no") & \linebreakcell{l}{$f \from \Sigma^{*} \to \Sigma^{*}$ \\ $L = \{(x, f(x))\}$} & \linebreakcell{l}{Prime factorization: \\ $2021 \rightarrow \{(2021, 43 \times 47)\}$} \\
\addlinespace[0.5em]
Search problem			& Answer is a relation of instance-solution pairs & \linebreakcell{l}{$R \subseteq \Sigma^{*} \times \Sigma^{*}$ \\ $A_{x} = \text{set of solutions for } x \in \Sigma^{*}$ \\ $L = \{(x,y) \mid (x,y) \in R, \ y \in A_{x}\}$} & \linebreakcell{l}{Nontrivial factors: \\ $6 \rightarrow \{(6,2), (6,3)\}$} \\
\addlinespace[0.5em]
Counting problem			& Answer is the number of solutions to a given search problem & \linebreakcell{l}{$R \subseteq \Sigma^{*} \times \Sigma^{*}$ \\ $A_{x} = \text{set of solutions for } x \in \Sigma^{*}$ \\ $L = \left| \{(x,y) \mid (x,y) \in R, \ y \in A_{x}\} \right|$} & \linebreakcell{l}{\# of nontrivial factors: \\ $16 \rightarrow 3$} \\
\addlinespace[0.5em]
Optimization problem		& Answer is the ``best possible" solution among a set of all feasible solutions & \linebreakcell{l}{Find $x \in \Sigma^{*}$ such that \\ $f(x)$ is minimal/maximal} & \linebreakcell{l}{Largest prime factor: \\ $93 \rightarrow 31$} \\
\bottomrule
\end{tabular}
\end{sidewaystable}

\section{Decidable Problems for Regular Languages}\label{sec:decidableREG}

\firstwords{We begin by considering} our common decision problems applied to models of computation that recognize the class of regular languages. Regular languages (and the associated models that recognize such languages) are very useful for practical applications since, as we will see, each of the common decision problems are decidable for this class. The downside, of course, is that the class of regular languages is much smaller than the other language classes we know, which in turn limits our expressive power.

\subsection*{Membership Problem}

The membership problem is a fundamental decision problem for any model of computation, since so much that pertains to decidability boils down to simply being able to figure out whether some model accepts some input word. Here, we will obtain our first decidability result by showing that there exists an algorithm that allows us to determine whether or not some deterministic finite automaton $\mathcal{B}$ accepts its input word $w$.

The technique we will apply in the proof of this result (and others) involves the use of a Turing machine to simulate the computation of the finite automaton. Then, whether the finite automaton accepts or doesn't accept the input word, the Turing machine returns the same result.

\begin{theorem}\label{thm:ADFAdecidable}
$\mathit{A}_{\DFA}$ is decidable.

\begin{proof}
Construct a Turing machine $\mathcal{M}_{\mathrm{ADFA}}$ that takes as input $\langle \mathcal{B}, w \rangle$, where $\mathcal{B}$ is a deterministic finite automaton and $w$ is the input word to $\mathcal{B}$, and performs the following steps:
\tmalgorithm{$\mathcal{M}_{\mathrm{ADFA}}$}{
\begin{enumerate}
\item Simulate $\mathcal{B}$ on input $w$.
\item If the simulation ends in an accepting state of $\mathcal{B}$, then accept. Otherwise, reject.
\end{enumerate}
}
This Turing machine accepts its input $\langle \mathcal{B}, w \rangle$ if and only if the deterministic finite automaton $\mathcal{B}$ accepts its input $w$, and the Turing machine rejects if and only if $\mathcal{B}$ rejects. Since the length of $w$, and therefore the simulation of $\mathcal{B}$, is finite, the Turing machine will never become trapped in an infinite computation. Thus, $\mathcal{M}_{\mathrm{ADFA}}$ decides the membership problem for deterministic finite automata.
\end{proof}
\end{theorem}

Since we know also that we can convert from nondeterministic finite automata to deterministic finite automata, and from regular expressions to deterministic finite automata, we get similar positive decidability results for these models.

\begin{corollary}
$\mathit{A}_{\NFA}$ and $\mathit{A}_{\RE}$ are decidable.

\begin{proof}
Given a nondeterministic finite automaton $\mathcal{B}$, we can convert it to an equivalent deterministic finite automaton $\mathcal{B}'$ and run $\mathcal{M}_{\mathrm{ADFA}}$ from the proof of Theorem~\ref{thm:ADFAdecidable} on the input $\langle \mathcal{B}', w \rangle$.

Likewise, given a regular expression $R$, we can convert it to an equivalent deterministic finite automaton $\mathcal{S}$ and run $\mathcal{M}_{\mathrm{ADFA}}$ on the input $\langle \mathcal{S}, w \rangle$.
\end{proof}
\end{corollary}

Indeed, one consequence of Kleene's theorem is that a decision problem being decidable for the class \DFA\ implies that it is also decidable for the classes \NFA\ and \RE\ (as well as the class \ENFA, though this arguably already falls under the ``nondeterministic" umbrella). Thus, going forward in this section, we will only focus on full proofs for the class \DFA.

\subsection*{Emptiness Problem}

Let's now turn to the emptiness problem. What does it mean for the language of a finite automaton to be empty? Obviously, it means the finite automaton doesn't accept any input words, but under what condition is this the case? Since any accepting computation of a finite automaton must conclude in a final state, we can reason that a finite automaton is incapable of accepting any input words only if there exists no path from its initial state to a final state.

We will use this reasoning in the algorithm to decide the emptiness problem for deterministic finite automata. Since every finite automaton has a finite set of states, we can traverse the transitions of the finite automaton starting from the initial state and mark each state as we encounter it. If, by the end of this traversal, we never mark a final state, then this implies there exists no path from the initial state to any final state.

\begin{theorem}\label{thm:EDFAdecidable}
$\mathit{E}_{\DFA}$ is decidable.

\begin{proof}
Construct a Turing machine $\mathcal{M}_{\mathrm{EDFA}}$ that takes as input $\langle \mathcal{B} \rangle$, where $\mathcal{B}$ is a deterministic finite automaton, and performs the following steps:
\tmalgorithm{$\mathcal{M}_{\mathrm{EDFA}}$}{
\begin{enumerate}
\item Mark the initial state of $\mathcal{B}$.
\item Mark all states that have an incoming transition from any previously marked state. Repeat until no new states are marked.
\item If no final state of $\mathcal{B}$ is marked, then accept. Otherwise, reject.
\end{enumerate}
}
This Turing machine accepts its input $\langle \mathcal{B} \rangle$ if and only if the deterministic finite automaton $\mathcal{B}$ has no possible computation path leading from its initial state to a final state, and the Turing machine rejects if and only if an accepting computation path exists in $\mathcal{B}$. Since the number of states of $\mathcal{B}$ that could be marked is finite, the Turing machine will never become trapped in an infinite computation. Thus, $\mathcal{M}_{\mathrm{EDFA}}$ decides the emptiness problem for deterministic finite automata.
\end{proof}
\end{theorem}

Naturally, following the same lines of reasoning we developed for the membership problem, we can establish the following corollary.

\begin{corollary}
$\mathit{E}_{\NFA}$ and $\mathit{E}_{\RE}$ are decidable.
\end{corollary}

\subsection*{Universality Problem}

The universality problem goes hand in hand with the emptiness problem we just studied. Indeed, the two problems are complementary: if $L = \Sigma^{*}$, then $\overline{L} = \emptyset$, where $\overline{L}$ denotes the complement of the language $L$ as we saw before. Intuitively, this makes sense: if a finite automaton accepts \emph{every} input word, then the complement of that finite automaton must accept \emph{no} input words.

Because of this complementary relationship, all we need to do in order to decide the universality problem is complement the language of our given deterministic finite automaton. Then, we can simply reuse the Turing machine $\mathcal{M}_{\mathrm{EDFA}}$ and the algorithm that we developed to decide the emptiness problem on the complemented finite automaton. In other words, in order to decide whether $L(\mathcal{B}) = \Sigma^{*}$, we simply need to decide whether $\overline{L(\mathcal{B})} = \emptyset$.

\begin{theorem}\label{thm:UDFAdecidable}
$\mathit{U}_{\DFA}$ is decidable.

\begin{proof}
Construct a Turing machine $\mathcal{M}_{\mathrm{UDFA}}$ that takes as input $\langle \mathcal{B} \rangle$, where $\mathcal{B}$ is a deterministic finite automaton, and performs the following steps:
\tmalgorithm{$\mathcal{M}_{\mathrm{UDFA}}$}{
\begin{enumerate}
\item Convert $\mathcal{B}$ to a deterministic finite automaton $\mathcal{B}'$ recognizing the language $\overline{L(\mathcal{B})}$ using the construction from the proof of Theorem~\ref{thm:FAclosurecomplement}.
\item Run $\mathcal{M}_{\mathrm{EDFA}}$ from the proof of Theorem~\ref{thm:EDFAdecidable} on input $\langle \mathcal{B}' \rangle$.
\item If $\mathcal{M}_{\mathrm{EDFA}}$ accepts, then accept. Otherwise, reject.
\end{enumerate}
}
This Turing machine accepts its input $\langle \mathcal{B} \rangle$ if and only if $\mathcal{M}_{\mathrm{EDFA}}$ accepts its input $\langle \mathcal{B}' \rangle$, where $\mathcal{B}'$ is the complement of the deterministic finite automaton $\mathcal{B}$. This occurs only when the language of $\mathcal{B}'$ is empty, and thus $\mathcal{B}$ must accept all input words in this case. Likewise, the Turing machine rejects if and only if $\mathcal{M}_{\mathrm{EDFA}}$ rejects, which means that the language of $\mathcal{B}'$ is not empty and there exists at least one input word not accepted by $\mathcal{B}$. Since both the conversion process to obtain $\mathcal{B}'$ as well as the process of running $\mathcal{M}_{\mathrm{EDFA}}$ are finite, the Turing machine will never become trapped in an infinite computation. Thus, $\mathcal{M}_{\mathrm{UDFA}}$ decides the universality problem for deterministic finite automata.
\end{proof}
\end{theorem}

Additionally, the following corollary holds as we would expect.

\begin{corollary}
$\mathit{U}_{\NFA}$ and $\mathit{U}_{\RE}$ are decidable.
\end{corollary}

\subsection*{Equivalence Problem}

Given two deterministic finite automata $\mathcal{B}$ and $\mathcal{C}$, how can we test whether $L(\mathcal{B}) = L(\mathcal{C})$? In theory, we could give all the words in $L(\mathcal{B})$ as input to $\mathcal{C}$ and test their membership using our Turing machine $\mathcal{M}_{\mathrm{ADFA}}$, and vice versa. However, this won't work out very well if either of $L(\mathcal{B})$ or $L(\mathcal{C})$ is infinite.

Instead, we can make the following useful observation: if $L(\mathcal{B}) = L(\mathcal{C})$, then every word in $L(\mathcal{B})$ and every word in $L(\mathcal{C})$ must appear in the intersection of the two languages. Equivalence means that no word belongs only to one of the two languages, and so in order to test whether two languages are equivalent, we need only test whether the non-intersecting parts of each language are empty: another application for our emptiness-testing Turing machine $\mathcal{M}_{\mathrm{EDFA}}$!

\begin{figure}
\centering
\begin{tikzpicture}[filled/.style={draw=black, fill=\fourthcolour, thick},outline/.style={draw=black, thick}]
	\draw[filled, even odd rule] (0,0) circle (1.5cm) node[xshift=-0.5cm] {$L(\mathcal{B})$}
							(0:2cm) circle (1.5cm) node[xshift=0.5cm] {$L(\mathcal{C})$};
\end{tikzpicture}
\caption{The symmetric difference of two languages $L(\mathcal{B})$ and $L(\mathcal{C})$}
\label{fig:symmetricdifference}
\end{figure}

The ``non-intersecting part" of two languages is more properly referred to as the \emph{symmetric difference} of the languages. Given two languages $L(\mathcal{B})$ and $L(\mathcal{C})$, their symmetric difference is the language
\begin{equation*}
L(\mathcal{B}) \symdiff L(\mathcal{C}) = \left( L(\mathcal{B}) \cap \overline{L(\mathcal{C})} \right) \cup \left( \overline{L(\mathcal{B})} \cap L(\mathcal{C}) \right).
\end{equation*}
The symmetric difference of $L(\mathcal{B})$ and $L(\mathcal{C})$ is illustrated by the shaded regions of the Venn diagram shown in Figure~\ref{fig:symmetricdifference}.

Since we know from earlier that the class of regular languages is closed under union, complement, and intersection, we can combine all of these closure properties to construct a deterministic finite automaton $\mathcal{D}$ whose language is $L(\mathcal{D}) = L(\mathcal{B}) \symdiff L(\mathcal{C})$. We will use this finite automaton $\mathcal{D}$ in our decision algorithm for the equivalence problem.

The idea behind this decision algorithm, as we mentioned before, is to test equivalence by testing the emptiness of the symmetric difference language. If $L(\mathcal{D})$ is empty, then we know that all words belong to the intersection of $L(\mathcal{B})$ and $L(\mathcal{C})$, and therefore $L(\mathcal{B}) = L(\mathcal{C})$.

\begin{theorem}\label{thm:EQDFAdecidable}
$\mathit{EQ}_{\DFA}$ is decidable.

\begin{proof}
Construct a Turing machine $\mathcal{M}_{\mathrm{EQDFA}}$ that takes as input $\langle \mathcal{B}, \mathcal{C} \rangle$, where $\mathcal{B}$ and $\mathcal{C}$ are deterministic finite automata, and performs the following steps:
\tmalgorithm{$\mathcal{M}_{\mathrm{EQDFA}}$}{
\begin{enumerate}
\item Construct a deterministic finite automaton $\mathcal{D}$ recognizing the language $L(\mathcal{D}) = L(\mathcal{B}) \symdiff L(\mathcal{C})$.
\item Run $\mathcal{M}_{\mathrm{EDFA}}$ from the proof of Theorem~\ref{thm:EDFAdecidable} on input $\langle \mathcal{D} \rangle$.
\item If $\mathcal{M}_{\mathrm{EDFA}}$ accepts, then accept. Otherwise, reject.
\end{enumerate}
}
This Turing machine accepts its input $\langle \mathcal{B}, \mathcal{C} \rangle$ if and only if $\mathcal{M}_{\mathrm{EDFA}}$ accepts its input $\langle \mathcal{D} \rangle$, where $\mathcal{D}$ is the deterministic finite automaton recognizing the symmetric difference language. This occurs only when the symmetric difference language is empty, and thus all of the words in both $L(\mathcal{B})$ and $L(\mathcal{C})$ appear in the intersection of these two languages; that is, $L(\mathcal{B}) = L(\mathcal{C})$. Likewise, the Turing machine rejects if and only if $\mathcal{M}_{\mathrm{EDFA}}$ rejects, which means that the symmetric difference language is nonempty and so $L(\mathcal{B}) \neq L(\mathcal{C})$. Since both the process of constructing $\mathcal{D}$ and the process of running $\mathcal{M}_{\mathrm{EDFA}}$ are finite, the Turing machine will never become trapped in an infinite computation. Thus, $\mathcal{M}_{\mathrm{EQDFA}}$ decides the equivalence problem for deterministic finite automata.
\end{proof}
\end{theorem}

Of course, we immediately have the usual corollary to accompany this result.

\begin{corollary}
$\mathit{EQ}_{\NFA}$ and $\mathit{EQ}_{\RE}$ are decidable.
\end{corollary}

\subsection*{Inclusion Problem}

\begin{construction}
Soon enough, we'll have a positive decidability result for the inclusion problem to go along with our collection of other positive decidability results.
\end{construction}
\section{Decidable Problems for Context-Free Languages}\label{sec:decidableCFL}

\firstwords{Moving on to the class} of context-free languages, we will again consider each of the common decision problems in turn, but here we must make a decision: much like we chose to prove regular language decidability results with the deterministic finite automaton model, do we prove context-free decidability results with the pushdown automaton model or with the context-free grammar model?

If we were to go with the pushdown automaton approach, we would have the familiarity in our proofs of simulating the computation of a simpler machine on a Turing machine, exactly as we did in the previous section. On the other hand, we would need to devise some way of managing the stack of this machine as we simulate its computation, and this added complexity would make our proofs somewhat more difficult.

By contrast, going with the context-free grammar approach, we can easily simulate a derivation (which is a computation in disguise) just by manipulating the rules of the context-free grammar in an appropriate manner. Since we know that context-free grammars and pushdown automata are equivalent in terms of recognition power, we lose nothing by selecting the grammar model over the automaton model, and so we will make our lives easier and take the grammar approach in this section.

\subsection*{Membership Problem}

As before, we begin with the simplest of decision problems. If we're given an input $\langle G, w \rangle$, where $G$ is a context-free grammar and $w$ is a word over the grammar's terminal alphabet $\Sigma$, the membership problem asks whether $G$ is capable of generating the word $w$.

The na\"{\i}ve approach to determining whether $G$ generates $w$ would have us check every possible derivation using the rules of $G$. However, this isn't a good approach, since in the worst case our algorithm may need to check infinitely many derivations. In fact, if we tried this approach and $G$ actually couldn't generate $w$, then our algorithm would keep checking derivations fruitlessly and would never halt! Thus, while the na\"{\i}ve approach \emph{semidecides} the problem, it doesn't \emph{decide} the problem.

We must somehow ensure that we only check some finite number of derivations in the process of testing membership. Thinking back to our discussion of Chomsky normal form, we made one important observation that will help us greatly: recall that if a context-free grammar in Chomsky normal form generates a word $w$, then the derivation of $w$ will take $2|w| - 1$ steps. This allows us to place an upper bound on the length of the derivation and, therefore, a limit on the number of steps performed by our membership-testing algorithm!

Using this observation, our decision algorithm will first convert its given context-free grammar $G$ to Chomsky normal form and then only check candidate derivations of $w$ that take $2|w| - 1$ steps. If $w$ can indeed be generated by $G$, then its derivation will be found by the algorithm.

\begin{theorem}\label{thm:ACFGdecidable}
$\mathit{A}_{\CFG}$ is decidable.

\begin{proof}
Construct a Turing machine $\mathcal{M}_{\mathrm{ACFG}}$ that takes as input $\langle G, w \rangle$, where $G$ is a context-free grammar and $w$ is a word, and performs the following steps:
\tmalgorithm{$\mathcal{M}_{\mathrm{ACFG}}$}{
\begin{enumerate}
\item Convert $G$ to an equivalent grammar $G'$ in Chomsky normal form.
\item Check the length of $w$:
	\begin{itemize}
	\item If $|w| = 0$, then list all derivations using $G'$ that take a single step.
	\item If $|w| \geq 1$, then list all derivations using $G'$ that take $2|w| - 1$ steps.
	\end{itemize}
\item If any of these derivations generate $w$, then accept. Otherwise, reject.
\end{enumerate}
}
This Turing machine accepts its input $\langle G, w \rangle$ if and only if the context-free grammar $G$ is capable of generating the word $w$; that is, if and only if $w \in L(G)$. Likewise, the Turing machine rejects if and only if no derivation exists that allows $G$ to generate $w$; that is, if and only if $w \not\in L(G)$. Since the process of converting $G$ to Chomsky normal form as well as the process of listing all derivations of a given length are finite, the Turing machine will never become trapped in an infinite computation. Thus, $\mathcal{M}_{\mathrm{ACFG}}$ decides the membership problem for context-free grammars.
\end{proof}
\end{theorem}

Unsurprisingly, the same positive decidability result holds for the pushdown automaton model.

\begin{corollary}
$\mathit{A}_{\PDA}$ is decidable.

\begin{proof}
Given a pushdown automaton $\mathcal{P}$, we can convert it to an equivalent context-free grammar $G$ and run $\mathcal{M}_{\mathrm{ACFG}}$ from the proof of Theorem~\ref{thm:ACFGdecidable} on the input $\langle G, w \rangle$.
\end{proof}
\end{corollary}

Reviewing our proof of Theorem~\ref{thm:ACFGdecidable} may make it more evident why we chose context-free grammars as our model in this section instead of pushdown automata: with grammars, we can obtain a concrete upper bound on how much work our Turing machine has to do before producing an answer.

On the other hand, if we're given a pushdown automaton as input, then we know we can simulate the computation of this automaton on a Turing machine in a rather straightforward manner. However, there's no way for us to predict in advance whether some branch of the pushdown automaton's computation tree will go on forever without accepting. Consider, for example, the branch of the computation tree that corresponds to following these two transitions repeatedly:
\begin{center}
\begin{tikzpicture}[node distance=3cm, >=latex, every state/.style={fill=white}]
\node[draw=none] (q) {$\dots$};
\node[state, minimum size=1cm] (q0) [right=1cm of q] {$q_{i}$};
\node[state, minimum size=1cm] (q1) [right of=q0] {$q_{i+1}$};
\node[draw=none] (qr) [right=1cm of q1] {$\dots$};

\path[-latex] (q) edge [above] node {} (q0);
\path[-latex] (q0) edge [bend left, above] node {$\epsilon, \epsilon \mapsto A$} (q1);
\path[-latex] (q1) edge [bend left, below] node {$\epsilon, A \mapsto \epsilon$} (q0);
\path[-latex] (q1) edge [above] node {} (qr);
\end{tikzpicture}
\end{center}
If our Turing machine na\"{\i}vely simulates the computation of the pushdown automaton by descending into individual branches of the computation tree one-by-one, then it could fall into this infinite-length branch trap. As a consequence, we can't guarantee that our Turing machine decides whether or not the pushdown automaton accepts its input word, and so we can't establish a positive decidability result for the membership problem unless we put in the additional work to avoid infinite-length branches in the computation tree.

By comparison, doesn't our context-free grammar approach seem so much more appealing?

\subsection*{Emptiness Problem}

For the emptiness problem, we again have a na\"{\i}ve approach to check whether $L(G) = \emptyset$ for some context-free grammar $G$: for all words $w$ over the terminal alphabet $\Sigma$, verify that there exists no derivation $S \Rightarrow^{*} w$ from the start nonterminal $S$ to some sequence of terminal symbols $w$. This is, of course, not a good approach for the same reason as before: the pesky notion of infinity gets in our way. There is an infinite number of words over $\Sigma$ for which we must test membership, so if the language $L(G)$ truly were empty, an algorithm following this approach would never halt. We would just keep checking and rejecting candidate words!

Instead, we will take an opposite approach to testing the emptiness of the grammar's language. Instead of checking that no sequence of \emph{terminal} symbols is generated by the grammar, we will check, for each \emph{nonterminal} symbol, whether that individual symbol yields some sequence of previously marked terminal symbols. If this is the case, then we will mark that nonterminal symbol. Then, whenever a marked nonterminal symbol appears in a derivation, we know that some sequence of terminal symbols will appear later in the derivation, and thus the grammar's language is nonempty.

Our decision algorithm for $E_{\CFG}$ works a lot like our decision algorithm for $E_{\DFA}$, except backwards: while the algorithm for $E_{\DFA}$ marked states starting from the initial state and leading to a final state, our algorithm for $E_{\CFG}$ will mark symbols starting from the terminal symbols and returning to the start nonterminal.

\begin{theorem}\label{thm:ECFGdecidable}
$\mathit{E}_{\CFG}$ is decidable.

\begin{proof}
Construct a Turing machine $\mathcal{M}_{\mathrm{ECFG}}$ that takes as input $\langle G \rangle$, where $G$ is a context-free grammar, and performs the following steps:
\tmalgorithm{$\mathcal{M}_{\mathrm{ECFG}}$}{
\begin{enumerate}
\item Mark all terminal symbols in $G$.
\item If $G$ has a rule of the form $A \rightarrow \alpha_{1} \dots \alpha_{k}$ and each symbol $\alpha_{1}, \dots, \alpha_{k}$ has already been marked, then mark the symbol $A$. Repeat until no new symbols are marked.
\item If the start nonterminal of $G$ has not been marked, then accept. Otherwise, reject.
\end{enumerate}
}
This Turing machine accepts its input $\langle G \rangle$ if and only if there exists no derivation $S \Rightarrow^{*} w$ in the context-free grammar $G$ from the start nonterminal $S$ to some sequence of terminal symbols $w$, and the Turing machine rejects if and only if such a derivation exists in $G$. Since the number of terminal symbols to mark and the number of rules of $G$ to check are both finite, the Turing machine will never become trapped in an infinite computation. Thus, $\mathcal{M}_{\mathrm{ECFG}}$ decides the emptiness problem for context-free grammars.
\end{proof}
\end{theorem}

As we have come to expect, the same result holds for pushdown automata.

\begin{corollary}
$\mathit{E}_{\PDA}$ is decidable.
\end{corollary}
%\section{Undecidable Problems for Context-Free Languages}\label{sec:undecidableproblemscontextfree}

\firstwords{Having established what it means} for a computation history to be invalid, and having showed that there exists a connection between the invalid computation histories of a Turing machine and our context-free models of computation, we have all we need to establish the undecidability of our remaining decision problems for context-free languages.

\subsection*{Universality Problem}

When we established the decidability of the universality problem for regular languages, we relied on Theorem~\ref{thm:FAclosurecomplement}, which showed that the class of regular languages was closed under complement. Taking the complement of a given finite automaton's language allowed us to use our existing decision procedure for $\mathit{E}_{\DFA}$ to decide $\mathit{U}_{\DFA}$, as shown in Theorem~\ref{thm:UDFAdecidable}.

For context-free languages, we know by Theorem~\ref{thm:ECFGdecidable} that $\mathit{E}_{\CFG}$ is decidable, so one may be tempted to try a similar approach to establish the supposed decidability of $\mathit{U}_{\CFG}$. However, Theorem~\ref{thm:CFLnonclosurecomplement} tells us that the class of context-free languages is \emph{not} closed under complement, and so this approach is a non-starter.

Instead, using a straightforward argument that relies on the connection between invalid computation histories and context-free languages, we can show that $\mathit{U}_{\CFG}$ is in fact \emph{undecidable}---the first such result we have for a model weaker than Turing machines!

\begin{theorem}\label{thm:UCFGundecidable}
$\mathit{U}_{\CFG}$ is undecidable.

\begin{proof}
Suppose that $\mathcal{M}$ is an arbitrary Turing machine. By Lemma~\ref{lem:invalidcomputationhistoryCFL}, we can construct a context-free grammar $G$ having the property that $L(G) = \Sigma^{*}$ if and only if $L(\mathcal{M}) = \emptyset$. This is because if $\mathcal{M}$ accepts no input words, then every input word given to $\mathcal{M}$ would produce an invalid computation history, and the aforementioned lemma tells us that the set of invalid computation histories of $\mathcal{M}$ is a context-free language.

However, if it were possible for us to check whether $L(G) = \Sigma^{*}$, then it would also become possible for us to decide whether $L(\mathcal{M}) = \emptyset$, and we know by Theorem~\ref{thm:ETMundecidable} that $\mathit{E}_{\TM}$ is undecidable. Therefore, $\mathit{U}_{\CFG}$ must also be undecidable.
\end{proof}
\end{theorem}

At this point, take a moment to reflect on the argument we just used to establish the undecidability of $\mathit{U}_{\CFG}$. Unlike our previous undecidability proofs using many-one reductions, we didn't need to construct any Turing machines here. Rather, we simply relied on a property that is true of any Turing machine: the set of invalid computation histories is a context-free language, and so we can construct a grammar for this language.

\begin{dangerous}
Note that nowhere did we say the language of our grammar $G$ \emph{is} $\Sigma^{*}$; all we said is that $L(G) = \Sigma^{*}$ \emph{if and only if} $L(\mathcal{M}) = \emptyset$. The language of $G$ depends on whatever arbitrary Turing machine $\mathcal{M}$ we are given, but if we had a method of deciding whether the language of $G$ was universal, then that same method could decide whether the language of this Turing machine $\mathcal{M}$ is empty.
\end{dangerous}

Naturally, since context-free grammars and pushdown automata are equivalent, this negative decidability result transfers over to our other context-free model of computation.

\begin{corollary}\label{cor:UPDAundecidable}
$\mathit{U}_{\PDA}$ is undecidable.
\end{corollary}

\subsection*{Equivalence Problem}

Recall that we can test the equivalence of two languages $L_{1}$ and $L_{2}$ by testing the emptiness of their symmetric difference, where the symmetric difference operation is defined in terms of union, intersection, and complement:
\begin{equation*}
L_{1} \symdiff L_{2} = \left( L_{1} \cap \overline{L_{2}} \right) \cup \left( \overline{L_{1}} \cap L_{2} \right).
\end{equation*}
Now, again, we know by Theorem~\ref{thm:ECFGdecidable} that we can test the emptiness of a context-free grammar, so we can test emptiness for two individual, separate context-free grammars just as easily. However, the class of context-free languages is not closed under intersection (by Theorem~\ref{thm:CFLnonclosureintersection}), nor is it closed under complement (by Theorem~\ref{thm:CFLnonclosurecomplement}). Thus, we have no way of testing the emptiness of the symmetric difference of two context-free grammars, and so it is hopeless for us to expect that we can decide the equivalence problem for context-free grammars.

But how do we formally prove the undecidability of $\mathit{EQ}_{\CFG}$? We need only fix one of our two context-free grammars to be ``special", which will then allow us to use what we know about $\mathit{U}_{\CFG}$ to arrive at our desired outcome.

\begin{theorem}\label{thm:EQCFGundecidable}
$\mathit{EQ}_{\CFG}$ is undecidable.

\begin{proof}
Let $G_{1}$ be an arbitrary context-free grammar, and let $G_{2}$ be a context-free grammar that generates all words over its alphabet $\Sigma$. If it were possible for us to check whether $L(G_{1}) = L(G_{2})$, then it would also become possible for us to decide whether $L(G_{1}) = \Sigma^{*}$, and we know by Theorem~\ref{thm:UCFGundecidable} that $\mathit{U}_{\CFG}$ is undecidable. Therefore, $\mathit{EQ}_{\CFG}$ must also be undecidable.
\end{proof}
\end{theorem}

In this proof, we have reduced in the sense that we made an instance of the decision problem $\mathit{EQ}_{\CFG}$ appear like an instance of the decision problem $\mathit{U}_{\CFG}$, and since the latter problem is undecidable, so too must the former problem be undecidable.

As we would expect, the same outcome holds for pushdown automata.

\begin{corollary}\label{cor:EQPDAundecidable}
$\mathit{EQ}_{\PDA}$ is undecidable.
\end{corollary}

\subsection*{Inclusion Problem}

\begin{construction}
Yes, the inclusion problem is also undecidable for context-free languages. This fact will be proved in due time.
\end{construction}
\section{An Undecidable Problem for Turing Machines}\label{sec:undecidableTM}

\firstwords{After studying so many} decision problems with positive decidability results, one might become convinced that our models of computation are capable of deciding anything we throw at it. Unfortunately, this is not the case: our regular and context-free models happen to be just simple enough to allow us to obtain the positive decidability results we saw in the previous sections, but turning to our strongest model of computation---the Turing machine---things begin to come apart.

But wait, since Turing machines are so powerful, doesn't that mean they can solve all kinds of problems? Indeed, that's true, but they can't \emph{decide} many problems. Remember, in order to decide a problem, the Turing machine must always halt and either accept or reject the input word it was given. As we observed earlier, the most common issue the Turing machine may encounter comes in the ``halt" step, since there's no guarantee that the machine will halt on every input word.

Indeed, even the most basic of decision problems is rendered undecidable on a Turing machine, simply because of that fundamental limitation that the machine may get caught in an infinite loop during its computation. We've proved that the membership problems for regular languages and for context-free languages are both decidable by virtue of the fact that the models recognizing such classes of languages always halt and either accept and reject their input words. For Turing machines, though, we lose this valuable decidability property.

Before we continue, let's review the notions of decidability and semidecidability from the previous chapter. Suppose that $\mathcal{M}$ is a Turing machine recognizing a language $L(\mathcal{M})$. Recall that:
\begin{colouredbox}
\begin{itemize}
\item $L(\mathcal{M})$ is \emph{decidable} if
	\begin{itemize}
	\item whenever $w \in L(\mathcal{M})$, $\mathcal{M}$ accepts $w$, while
	\item whenever $w \not\in L(\mathcal{M})$, $\mathcal{M}$ rejects $w$; and
	\end{itemize}
\item $L(\mathcal{M})$ is \emph{semidecidable} if
	\begin{itemize}
	\item whenever $w \in L(\mathcal{M})$, $\mathcal{M}$ accepts $w$, while
	\item whenever $w \not\in L(\mathcal{M})$, $\mathcal{M}$ either rejects $w$ or enters an infinite loop.
	\end{itemize}
\end{itemize}
\end{colouredbox}

It's quite straightforward to show that the membership problem for Turing machines, $\mathit{A}_{\TM}$, is at least semidecidable. All we need to do is simulate the computation of $\mathcal{M}$ on $w$! No matter what $\mathcal{M}$ does---accept, reject, or loop forever---its behaviour matches our definition of semidecidability.

\begin{theorem}\label{thm:ATMsemidecidable}
$\mathit{A}_{\TM}$ is semidecidable.

\begin{proof}
Construct a Turing machine $\mathcal{M}_{\mathrm{ATM}}$ that takes as input $\langle \mathcal{M}, w \rangle$, where $\mathcal{M}$ is a Turing machine and $w$ is a word, and performs the following steps:
\tmalgorithm{$\mathcal{M}_{\mathrm{ATM}}$}{
\begin{enumerate}
\item Simulate $\mathcal{M}$ on input $w$.
\item If $\mathcal{M}$ ever enters its accepting state $q_{\text{accept}}$, then accept. If $\mathcal{M}$ ever enters its rejecting state $q_{\text{reject}}$, then reject.
\end{enumerate}
}
This Turing machine accepts its input $\langle \mathcal{M}, w \rangle$ if and only if the Turing machine $\mathcal{M}$ accepts its input $w$, and the Turing machine rejects if and only if $\mathcal{M}$ rejects. If $\mathcal{M}$ enters an infinite loop on the input $w$, then $\mathcal{M}_{\mathrm{ATM}}$ will likewise loop forever. Thus, $\mathcal{M}_{\mathrm{ATM}}$ semidecides the membership problem for Turing machines.
\end{proof}
\end{theorem}

The machine $\mathcal{M}_{\mathrm{ATM}}$ that we constructed in the proof of Theorem~\ref{thm:ATMsemidecidable} looks quite similar to the machine $\mathcal{M}_{\mathrm{ADFA}}$ we constructed in the proof of Theorem~\ref{thm:ADFAdecidable} to decide $\mathit{A}_{\DFA}$. However, unlike the machine $\mathcal{M}_{\mathrm{ADFA}}$, this machine $\mathcal{M}_{\mathrm{ATM}}$ only semidecides $\mathit{A}_{\TM}$, since it has no way of explicitly rejecting its input word if it gets caught in an infinite loop. The machine $\mathcal{M}_{\mathrm{ATM}}$ can't even guess that it's looping infinitely and ``short circuit" its computation to automatically reject, for the plain reason that if the machine has performed $1\,000\,000$ steps, say, then there's no general way of determining whether it will reach the accepting or rejecting state on step $1\,000\,001$.

Note also that $\mathcal{M}_{\mathrm{ATM}}$ is the canonical example of a universal Turing machine, since we can give it an encoding of any Turing machine $\mathcal{M}$ and any input word $w$ and it will simulate the computation of that Turing machine on that input word.

Now, getting back to our main point, we must prove that $\mathit{A}_{\TM}$ is undecidable; that is to say, $\mathit{A}_{\TM}$ is \emph{not decidable}. The technique we will use is essentially a proof by contradiction: we will assume that we have some machine capable of deciding $\mathit{A}_{\TM}$, and then show that the existence of such a machine leads to a logical absurdity when we give a nefariously crafted input word to that machine.

\begin{theorem}\label{thm:ATMundecidable}
$\mathit{A}_{\TM}$ is undecidable.

\begin{proof}
Assume by way of contradiction that $\mathit{A}_{\TM}$ is decidable, and suppose that $\mathcal{M}_{\mathrm{A}}$ is a Turing machine that decides $\mathit{A}_{\TM}$. Given an encoding of a Turing machine $\mathcal{M}$ and an input word $w$, $\mathcal{M}_{\mathrm{A}}$ produces the same answer that $\mathcal{M}$ would produce:
\begin{itemize}
\item $\mathcal{M}_{\mathrm{A}}$ accepts $\langle \mathcal{M}, w \rangle$ if $\mathcal{M}$ accepts $w$; and
\item $\mathcal{M}_{\mathrm{A}}$ rejects $\langle \mathcal{M}, w \rangle$ if $\mathcal{M}$ rejects $w$.
\end{itemize}
Note that we don't know exactly how $\mathcal{M}_{\mathrm{A}}$ decides the membership problem: we only assume that it is capable of doing it.

We now construct a new Turing machine $\mathcal{X}$ that relies on using $\mathcal{M}_{\mathrm{A}}$ as an intermediate step of its computation. The machine $\mathcal{X}$ takes as input $\langle \mathcal{M} \rangle$, where $\mathcal{M}$ is a Turing machine, and performs the following steps:
\tmalgorithm{$\mathcal{X}$}{
\begin{enumerate}
\item Run $\mathcal{M}_{\mathrm{A}}$ on input $\langle \mathcal{M}, \langle \mathcal{M} \rangle \rangle$.
\item If $\mathcal{M}_{\mathrm{A}}$ accepts, then \emph{reject}. If $\mathcal{M}_{\mathrm{A}}$ rejects, then \emph{accept}.
\end{enumerate}
}
\noindent
Here, we see that $\mathcal{X}$ gives as input to $\mathcal{M}_{\mathrm{A}}$ exactly what $\mathcal{M}_{\mathrm{A}}$ expects to receive as input: an encoding of a Turing machine $\mathcal{M}$, together with an input word to $\mathcal{M}$ that happens to be the encoded description of the Turing machine $\mathcal{M}$ itself. This is allowed, since an encoded description of a Turing machine is nothing more than a sequence of symbols just like any other input word. After $\mathcal{M}_{\mathrm{A}}$ halts and produces its answer, $\mathcal{X}$ produces the opposite answer to whatever $\mathcal{M}_{\mathrm{A}}$ gave:
\begin{itemize}
\item $\mathcal{X}$ accepts $\langle \mathcal{M} \rangle$ if $\mathcal{M}$ rejects $\langle \mathcal{M} \rangle$; and
\item $\mathcal{X}$ rejects $\langle \mathcal{M} \rangle$ if $\mathcal{M}$ accepts $\langle \mathcal{M} \rangle$.
\end{itemize}
Figure~\ref{fig:ATMproofD} depicts the behaviour of $\mathcal{X}$ visually. Observe in this figure that $\mathcal{M}_{\mathrm{A}}$ is a ``black box", because as we remarked earlier, we don't know exactly how $\mathcal{M}_{\mathrm{A}}$ decides $\mathit{A}_{\TM}$.

Now, consider what happens when we give $\mathcal{X}$ an encoding of \emph{itself} as input. Suppose we give $\mathcal{X}$ the input $\langle \mathcal{X} \rangle$, so that $\mathcal{X}$ runs $\mathcal{M}_{\mathrm{A}}$ on the input $\langle \mathcal{X}, \langle \mathcal{X} \rangle \rangle$. Tracing through the behaviour of $\mathcal{M}_{\mathrm{A}}$ followed by $\mathcal{X}$, we can see that:
\begin{itemize}
\item $\mathcal{X}$ accepts $\langle \mathcal{X} \rangle$ if $\mathcal{X}$ rejects $\langle \mathcal{X} \rangle$; and
\item $\mathcal{X}$ rejects $\langle \mathcal{X} \rangle$ if $\mathcal{X}$ accepts $\langle \mathcal{X} \rangle$.
\end{itemize}
In either case, the actual computation of $\mathcal{X}$ on the input $\langle \mathcal{X} \rangle$ must return the opposite answer as the simulated computation of $\mathcal{X}$ on the input $\langle \mathcal{X} \rangle$, which is of course impossible! As a result, we obtain a contradiction. Since the only assumption in our proof is about the existence of $\mathcal{M}_{\mathrm{A}}$, we conclude that no such machine $\mathcal{M}_{\mathrm{A}}$ can exist to decide $\mathit{A}_{\TM}$.
\end{proof}
\end{theorem}

\begin{figure}
\centering
\begin{tikzpicture}
\draw[draw=\fourthcolour, thick, fill=\fourthcolour, rounded corners] (1.5,-0.55) rectangle (9,4);
\draw[draw=\fifthcolour, thick, fill=\fifthcolour] (2,1) rectangle (3.5,2.5);
\draw[draw=black, thick, fill=black] (5,1) rectangle (6.5,2.5);

\node[draw=none, color=\maincolour] at (2.75,3.25) {\huge$\mathcal{X}$};
\node[draw=none, color=\maincolour, text width=2cm, align=center] at (2.75,1.75) {Convert \\ input};
\node[draw=none, color=white, text width=2cm, align=center] at (5.75,1.75) {\huge$\mathcal{M}_{\mathrm{A}}$};

\coordinate (A) at (2,1.75);
\coordinate (B) at (0.5,1.75);
\draw[-latex, color=\maincolour, thick] (B) -- (A);
\node[draw=none, color=\maincolour] at (0.95,2.15) {\footnotesize$\langle \mathcal{M} \rangle$};

\coordinate (C) at (5,1.75);
\coordinate (D) at (3.5,1.75);
\draw[-latex, color=\maincolour, thick] (D) -- (C);
\node[draw=none, color=\maincolour] at (4.25,2.15) {\footnotesize$\langle \mathcal{M}, \langle \mathcal{M} \rangle \rangle$};

\coordinate (E) at (11.3,3);
\coordinate (EE) at (11.3,0.5);
\coordinate (EF) at (6.5, 3);
\coordinate (F) at (5.75,2.5);
\coordinate (FE) at (6.5, 0.5);
\coordinate (FF) at (5.75,1);
\draw[-latex, color=\maincolour, thick] (F) -- (EF) -- (E);
\draw[-latex, color=\maincolour, thick] (FF) -- (FE) -- (EE);
\node[draw=none, color=\maincolour] at (7.25,3.3) {\footnotesize$\mathcal{M}_{\mathrm{A}}$ accepts $\langle \mathcal{M}, \langle \mathcal{M} \rangle \rangle$};
\node[draw=none, color=\maincolour] at (7.25,3.65) {\footnotesize$\mathcal{M}$ accepts $\langle \mathcal{M} \rangle$};
\node[draw=none, color=\maincolour] at (7.25,-0.2) {\footnotesize$\mathcal{M}_{\mathrm{A}}$ rejects $\langle \mathcal{M}, \langle \mathcal{M} \rangle \rangle$};
\node[draw=none, color=\maincolour] at (7.25,0.15) {\footnotesize$\mathcal{M}$ rejects $\langle \mathcal{M} \rangle$};
\node[draw=none, color=\maincolour] at (10.25,3.45) {\footnotesize$\mathcal{X}$ rejects $\langle \mathcal{M} \rangle$};
\node[draw=none, color=\maincolour] at (10.25,0) {\footnotesize$\mathcal{X}$ accepts $\langle \mathcal{M} \rangle$};
\end{tikzpicture}
\caption{The Turing machine $\mathcal{X}$ from the proof of Theorem~\ref{thm:ATMundecidable}}
\label{fig:ATMproofD}
\end{figure}

\subsubsection*{Diagonalization}

The underlying concept that makes the proof of Theorem~\ref{thm:ATMundecidable} work is \emph{Cantor's diagonal argument}, named for the mathematician Georg Cantor, who used the argument to demonstrate the existence of uncountably infinite sets~\citeyearpar{Cantor1891Mannigfaltigkeitslehre}. An \emph{uncountably infinite} set is one whose elements cannot be placed into a bijection with the \emph{countably infinite} set of natural numbers $\mathbb{N}$.

The crux of the diagonal argument is as follows: if we take a set $T$ of all infinite-length binary words and assume that this set is countably infinite, then we can enumerate (that is, \emph{count}) all of the elements of $T$, and we can denote this by writing something like $T = \{t_{0}, t_{1}, t_{2}, \dots\}$. However, using the set $T$ itself, we can construct a new word $s$ that doesn't belong to $T$ by taking the $i$th symbol of $s$ to be the complement of the $i$th symbol in the word $t_{i}$. This means that if the $i$th symbol in the word $t_{i}$ is \texttt{0}, then we set the $i$th symbol of $s$ to be \texttt{1}, and vice versa. This construction process is illustrated in Figure~\ref{fig:cantorsdiagonalargument}.

In the process of constructing this new word $s$, we can see that it differs from every other word in $T$ in at least one position; namely, $s$ differs from $t_{i}$ at least in position $i$. (For example, compare $s$ to $t_{9}$ in Figure~\ref{fig:cantorsdiagonalargument}.) Therefore, it's impossible for $s$ to have already been included in $T$. Since $s$ is not an element of $T$, we could not have accounted for $s$ in our enumeration of $T$. Moreover, since we can create an infinite number of new words in this way, and none of these new words is accounted for in our enumeration, $T$ must not be countably infinite.

\begin{figure}
\centering
\begin{tabular}{r c p{0.05em} p{0.05em} p{0.05em} p{0.05em} p{0.05em} p{0.05em} p{0.05em} p{0.05em} p{0.05em} p{0.05em} c}
$t_{0}$	& $=$	& {\leavevmode\color{\secondcolour}\ul{0}}	& 0	& 0	& 0	& 0	& 0	& 0	& 0	& 0	& 0	& $\cdots$ \\
$t_{1}$	& $=$	& 1	& {\leavevmode\color{\secondcolour}\ul{1}}	& 1	& 1	& 1	& 1	& 1	& 1	& 1	& 1	& $\cdots$ \\
$t_{2}$	& $=$	& 0	& 0	& {\leavevmode\color{\secondcolour}\ul{1}}	& 0	& 1	& 1	& 0	& 1	& 1	& 1	& $\cdots$ \\
$t_{3}$	& $=$	& 0	& 1	& 1	& {\leavevmode\color{\secondcolour}\ul{0}}	& 1	& 0	& 0	& 1	& 1	& 0	& $\cdots$ \\
$t_{4}$	& $=$	& 1	& 0	& 1	& 0	& {\leavevmode\color{\secondcolour}\ul{1}}	& 0	& 1	& 0	& 1	& 0	& $\cdots$ \\
$t_{5}$	& $=$	& 1	& 1	& 0	& 0	& 1	& {\leavevmode\color{\secondcolour}\ul{1}}	& 0	& 0	& 1	& 1	& $\cdots$ \\
$t_{6}$	& $=$	& 0	& 0	& 0	& 1	& 1	& 1	& {\leavevmode\color{\secondcolour}\ul{0}}	& 0	& 0	& 1	& $\cdots$ \\
$t_{7}$	& $=$	& 1	& 0	& 0	& 1	& 0	& 1	& 1	& {\leavevmode\color{\secondcolour}\ul{0}}	& 0	& 1	& $\cdots$ \\
$t_{8}$	& $=$	& 0	& 1	& 0	& 1	& 0	& 1	& 0	& 1	& {\leavevmode\color{\secondcolour}\ul{0}}	& 1	& $\cdots$ \\
$t_{9}$	& $=$	& 1	& 0	& 0	& 1	& 0	& 0	& 1	& 1	& 1	& {\leavevmode\color{\secondcolour}\ul{0}}	& $\cdots$ \\
		& $\vdots$	&	&	&	&	&	&	&	&	&	&	& \\
\midrule
$s$		& $=$	& {\leavevmode\color{\secondcolour}\ul{1}}	& {\leavevmode\color{\secondcolour}\ul{0}}	& {\leavevmode\color{\secondcolour}\ul{0}}	& {\leavevmode\color{\secondcolour}\ul{1}}	& {\leavevmode\color{\secondcolour}\ul{0}}	& {\leavevmode\color{\secondcolour}\ul{0}}	& {\leavevmode\color{\secondcolour}\ul{1}}	& {\leavevmode\color{\secondcolour}\ul{1}}	& {\leavevmode\color{\secondcolour}\ul{1}}	& {\leavevmode\color{\secondcolour}\ul{1}}	& $\cdots$
\end{tabular}
\caption{Constructing a new word $s$ that doesn't belong to the set $T$ using Cantor's diagonal argument}
\label{fig:cantorsdiagonalargument}
\end{figure}

Cantor's diagonal argument is popularly used to prove that the set of real numbers $\mathbb{R}$ is uncountably infinite, but the same argument can be applied to Turing machines in order to prove that $\mathit{A}_{\TM}$ is undecidable. If we could enumerate all Turing machines, then we could construct a table where each row corresponds to a Turing machine $\mathcal{M}_{i}$ and each column corresponds to an encoding of a Turing machine $\langle \mathcal{M}_{j} \rangle$. Then, the entries of the table could be taken to be the result of running our supposed machine $\mathcal{M}_{\mathrm{A}}$ on the input $\langle \mathcal{M}_{i}, \langle \mathcal{M}_{j} \rangle \rangle$, where each entry is either ``accept" or ``reject" depending on the answer that $\mathcal{M}_{\mathrm{A}}$ produced.

\begin{center}
\small
\begin{tabular}{c | c c c c c}
				& $\langle \mathcal{M}_{0} \rangle$	& $\langle \mathcal{M}_{1} \rangle$	& $\langle \mathcal{M}_{2} \rangle$	& $\langle \mathcal{M}_{3} \rangle$	& $\cdots$ \\
\hline
$\mathcal{M}_{0}$	& accept						& accept						& reject						& reject						& \\
$\mathcal{M}_{1}$	& accept						& reject						& reject						& accept						& \\
$\mathcal{M}_{2}$	& reject						& reject						& accept						& accept						& \\
$\mathcal{M}_{3}$	& accept						& accept						& accept						& reject						& \\
$\vdots$			&							&							&							&							& $\ddots$
\end{tabular}
\end{center}

By the fact that $\mathcal{X}$ itself is a Turing machine, we know that $\mathcal{X}$ must appear somewhere in our enumerated list, and thus in our table. We also know each entry in the row corresponding to $\mathcal{X}$: since $\mathcal{X}$ produces the opposite answer to whatever $\mathcal{M}_{\mathrm{A}}$ produces on the input $\langle \mathcal{M}_{i}, \langle \mathcal{M}_{i} \rangle \rangle$, the entry in column $i$ must be the opposite of the entry at position $(i,i)$ of the table. But again, therein lies the problem: what is the table entry at the position corresponding to $\langle \mathcal{X}, \langle \mathcal{X} \rangle \rangle$?

\begin{center}
\small
\begin{tabular}{c | c c c c c c c}
				& $\langle \mathcal{M}_{0} \rangle$	& $\langle \mathcal{M}_{1} \rangle$	& $\langle \mathcal{M}_{2} \rangle$	& $\langle \mathcal{M}_{3} \rangle$	& $\cdots$	& $\langle \mathcal{X} \rangle$	& $\cdots$ \\
\hline
$\mathcal{M}_{0}$	& {\leavevmode\color{\secondcolour}\ul{accept}}					& accept						& reject						& reject			& 					& accept	& \\
$\mathcal{M}_{1}$	& accept						& {\leavevmode\color{\secondcolour}\ul{reject}}						& reject						& accept			& 					& reject	& \\
$\mathcal{M}_{2}$	& reject						& reject						& {\leavevmode\color{\secondcolour}\ul{accept}}					& accept			& 					& reject	& \\
$\mathcal{M}_{3}$	& accept						& accept						& accept						& {\leavevmode\color{\secondcolour}\ul{reject}}			& 					& accept	& \\
$\vdots$			&							&							&							&							& $\ddots$		&					& \\
$\mathcal{X}$		& {\leavevmode\color{\secondcolour}\ul{reject}}		& {\leavevmode\color{\secondcolour}\ul{accept}}	& {\leavevmode\color{\secondcolour}\ul{reject}}		& {\leavevmode\color{\secondcolour}\ul{accept}}	&		& \textbf{?}	& \\
$\vdots$			&							&							&							&							&				&					& $\ddots$
\end{tabular}
\end{center}

The entry at this position of the table somehow needs to be the opposite of itself, and so we arrive at the same contradiction as before. The Turing machine $\mathcal{X}$ cannot exist, which means that the Turing machine $\mathcal{M}_{\mathrm{A}}$ cannot exist, which means that $\mathit{A}_{\TM}$ is undecidable.
\section{A Non-Semidecidable Problem for Turing Machines}\label{sec:nonsemidecidableTM}

\firstwords{Now that we know} the membership problem $\mathit{A}_{\TM}$ is semidecidable but not decidable, we can be fairly sure that obtaining a positive decidability result for other Turing machine decision problems is hopeless. After all, if we can't even decide the membership of a word in a Turing machine's language, what other questions about the language can we hope to decide?

In order to crush our spirits even further, let's prove another fact about languages and their relation to Turing machines. Earlier, we showed that the class of decidable languages was a proper subset of the class of semidecidable languages by proving that $\mathit{A}_{\TM}$ is semidecidable, but not decidable. The semidecidable languages are the ``largest" class we've studied thus far in this book. However, semidecidability is not the upper limit when it comes to language classes! We can in fact show via a familiar argument that there exist languages that lie even outside of the class of semidecidable languages.

\begin{theorem}\label{thm:nonsemidecidablelanguages}
There exist languages that are not semidecidable.

\begin{proof}
In order for a language to be (at least) semidecidable, it must be recognized by some Turing machine. Recall that we can describe any Turing machine $\mathcal{M}$ with a unique encoding $\langle \mathcal{M} \rangle$. Although there is an infinite number of Turing machines, each of these encodings has a finite length, so we can enumerate them (say, lexicographically). By associating each encoding with a natural number according to that encoding's position in our enumeration, we obtain a bijection between encodings and the set of natural numbers $\mathbb{N}$, which means that the set of all Turing machines is countably infinite.

Now, consider the set of all words over some alphabet $\Sigma$. Again, even though this set is infinite, each of these words has a finite length, so we can enumerate them (say, lexicographically); denote the enumerated set by $\{w_{0}, w_{1}, w_{2}, \dots\}$. From this, we can define languages in terms of the words $w_{i}$ that belong to the language: we represent each language as an infinite-length binary word where the $i$th symbol of the word is \texttt{1} if $w_{i} \in L$, and \texttt{0} otherwise.

However, the set of all languages cannot be countably infinite: by Cantor's diagonal argument, if we could enumerate all of the infinite-length binary words representing every language, then it would be possible to construct a new infinite-length word representing a language not included in our enumeration. Therefore, the set of all languages must be uncountably infinite, and so there must exist languages that are not recognized by any Turing machine.
\end{proof}
\end{theorem}

Looking a bit more closely at our proof, we find ourselves confronted by a perhaps-uncomfortable fact: not only do non-semidecidable languages exist, but \emph{uncountably many} non-semidecidable languages exist. Thus, if we were to choose a language uniformly at random from the set of all languages, that language will---in the probabilistic sense---almost never be decidable by any Turing machine. In other words, the probability that a computer can solve the problem corresponding to our randomly chosen language is \emph{zero}!

Theorem~\ref{thm:nonsemidecidablelanguages} is an example of a theorem having a \emph{non-constructive} proof: the proof tells us that there exist languages that are not semidecidable, but it doesn't actually give us an example of such a language. We will, of course, see such an example momentarily, but first we must come up with one more definition.

\subsubsection*{Co-Semidecidability}

The property of semidecidability, at its core, is non-symmetric. Recall the decision problem $\mathit{A}_{\TM}$: we are guaranteed to get an answer in the positive case, if the Turing machine accepts its input, but in the negative case if the Turing machine loops forever, we will never end up with an answer.

What if we wanted to guarantee answers in the negative case, though? Consider how we might define a ``looping" decision problem, where we want to decide whether a Turing machine \emph{will} loop forever on some input. In the positive case, if the Turing machine does in fact loop forever, we will again never end up with an answer. But in the negative case, if the Turing machine halts, we are guaranteed to get an answer! We can't reasonably call this new decision problem semidecidable, since its behaviour doesn't line up with our definition of semidecidability. We need some new terminology.

Let $\overline{L}$ denote the complement of a language $L$. If we know something about this language $L$---namely, that it is semidecidable---then we can draw a straightforward conclusion about the complement of this language.

\begin{definition}[Co-semidecidable language]\label{def:cosemidecidablelanguage}
If a language $L$ is semidecidable, then we say that the complement of this language, $\overline{L}$, is co-semidecidable. \medskip

\noindent
Alternatively, given a Turing machine $\mathcal{M}$, we say that the language $L(\mathcal{M})$ is co-semidecidable if,
\begin{itemize}
\item whenever $w \in L(\mathcal{M})$, then either $\mathcal{M}$ accepts $w$ or $\mathcal{M}$ enters an infinite loop; and
\item whenever $w \not\in L(\mathcal{M})$, then $\mathcal{M}$ rejects $w$.
\end{itemize}
\end{definition}

Observe that the definition of co-semidecidability takes the non-symmetric definition of semidecidability and flips it around: now, we are guaranteed an answer in the negative case when $w \not\in L(\mathcal{M})$, whereas semidecidability guaranteed us an answer in the positive case when $w \in L(\mathcal{M})$.

\begin{dangerous}
Saying that $\overline{L}$ is co-semidecidable is another way of saying that we can \emph{\ul{semidecide}} the \emph{\ul{co}}mplement of $\overline{L}$; that is, we can semidecide $\overline{\overline{L}} = L$.
\end{dangerous}

One might wonder, if we have a Turing machine that semidecides some language $L$, whether we need a separate Turing machine to co-semidecide $\overline{L}$? In fact, we don't always need a new machine: the original language $L$ and the language of all words \emph{not} in the complement language $\overline{L}$ are one and the same. (Think about why this is the case.) Thus, if we want to, we can use the existing Turing machine $\mathcal{M}$ that semidecides $L$ also to co-semidecide $\overline{L}$: if $\mathcal{M}$ accepts some input, then we know that input cannot belong to $\overline{L}$, and so we can treat it as if it were rejected.

Going a step further, if we know something about the co-semidecidability of a language, does this say anything about the semidecidability of the same language? By Definition~\ref{def:cosemidecidablelanguage}, we know that $\overline{L}$ is co-semidecidable whenever $L$ is semidecidable, but this in itself does not suggest or imply that $\overline{L}$ is semidecidable. If we could further show that $\overline{L}$ were also semidecidable, then $L$ would be co-semidecidable, again by our definition. In this situation where $L$ is both semidecidable and co-semidecidable, we can ``upgrade" the class to which $L$ belongs.

\begin{figure}
\centering
\begin{tikzpicture}
\draw[draw=\fourthcolour, thick, fill=\fourthcolour, rounded corners] (1,-1) rectangle (7.2,3.5);
\draw[draw=\fifthcolour, thick, fill=\fifthcolour] (2.75,1) rectangle (4.25,2.5);
\draw[draw=\fifthcolour, thick, fill=\fifthcolour] (2.75,-0.75) rectangle (4.25,0.75);

\node[draw=none, color=\maincolour] at (2,2.85) {\large$\mathcal{M}_{\text{D}}$};
\node[draw=none, color=\maincolour, text width=2cm, align=center] at (3.5,1.75) {\large$\mathcal{M}_{\text{SD}}$};
\node[draw=none, color=\maincolour, text width=2cm, align=center] at (3.5,0) {\large$\mathcal{M}_{\text{coSD}}$};

\coordinate (A) at (0.1,0.9);
\coordinate (B) at (1.5,0.9);
\coordinate (C) at (2,1.75);
\coordinate (D) at (2,0);
\coordinate (E) at (2.75,1.75);
\coordinate (F) at (2.75,0);
\draw[-latex, color=\maincolour, thick] (A) -- (B) -- (C) -- (E);
\draw[-latex, color=\maincolour, thick] (A) -- (B) -- (D) -- (F);
\node[draw=none, color=\maincolour] at (0.5,1.2) {\footnotesize$w$};

\coordinate (G) at (4.25,1.75);
\coordinate (H) at (4.25,0);
\coordinate (I) at (9.7,1.75);
\coordinate (J) at (9.7,0);
\draw[-latex, color=\maincolour, thick] (G) -- (I);
\draw[-latex, color=\maincolour, thick] (H) -- (J);
\node[draw=none, color=\maincolour] at (5.7,2.15) {\footnotesize$\mathcal{M}_{\text{SD}}$ accepts $w$};
\node[draw=none, color=\maincolour] at (5.7,0.35) {\footnotesize$\mathcal{M}_{\text{coSD}}$ accepts $w$};
\node[draw=none, color=\maincolour] at (8.5,2.15) {\footnotesize$\mathcal{M}_{\text{D}}$ accepts $w$};
\node[draw=none, color=\maincolour] at (8.5,0.35) {\footnotesize$\mathcal{M}_{\text{D}}$ rejects $w$};
\end{tikzpicture}
\caption{The Turing machine $\mathcal{M}_{\text{D}}$ from the proof of Theorem~\ref{thm:decidablesemidecidable}}
\label{fig:decidablesemidecidableproofMD}
\end{figure}

\begin{theorem}\label{thm:decidablesemidecidable}
A language $L$ is decidable if and only if $L$ is both semidecidable and co-semidecidable.

\begin{proof}
$(\Rightarrow)$: Suppose $L$ is decidable. Since all decidable languages are also semidecidable, we have that $L$ is semidecidable. If we take the Turing machine deciding $L$ and exchange the accepting and rejecting outputs, then we get a machine that decides $\overline{L}$. Thus, $\overline{L}$ is semidecidable, and so $L$ is co-semidecidable.

$(\Leftarrow)$: Suppose $L$ is both semidecidable and co-semidecidable. Then there exists a Turing machine $\mathcal{M}_{\text{SD}}$ semideciding $L$ and a Turing machine $\mathcal{M}_{\text{coSD}}$ co-semideciding $L$. Using these two Turing machines, we can construct a Turing machine $\mathcal{M}_{\text{D}}$ that takes as input a word $w$ and performs the following steps:
\tmalgorithm{$\mathcal{M}_{\text{D}}$}{
\begin{enumerate}
\item Run both $\mathcal{M}_{\text{SD}}$ and $\mathcal{M}_{\text{coSD}}$ on the input word $w$ in parallel.
\item If $\mathcal{M}_{\text{SD}}$ accepts, then accept. If $\mathcal{M}_{\text{coSD}}$ accepts, then reject.
\end{enumerate}
}
\noindent
Figure~\ref{fig:decidablesemidecidableproofMD} depicts the behaviour of $\mathcal{M}_{\text{D}}$ visually.

Since every word $w$ must belong to either $L$ or $\overline{L}$, either $\mathcal{M}_{\text{SD}}$ or $\mathcal{M}_{\text{coSD}}$ must accept $w$. Moreover, since $\mathcal{M}_{\text{D}}$ halts whenever either $\mathcal{M}_{\text{SD}}$ or $\mathcal{M}_{\text{coSD}}$ accepts, we have that $\mathcal{M}_{\text{D}}$ decides $L$, and so $L$ is decidable.
\end{proof}
\end{theorem}

\subsubsection*{The Non-membership Problem}

To motivate our study of co-semidecidability, we considered a ``looping" decision problem, but our formulation of this problem was rather ad hoc and imprecise. Since co-semidecidability is fundamentally a property of the complement of a language, let's further motivate our study by taking an existing decision problem that we have already formally defined and studying its complement. Namely, consider the Turing machine ``non-membership" problem:
\begin{align*}
\overline{\mathit{A}_{\TM}} = \{\langle \mathcal{M}, w \rangle \mid \ & \mathcal{M} \text{ is a Turing machine that} \\[-0.2em]
	&\text{ does \emph{not} accept input word } w\}.
\end{align*}

Taking all that we know about decidability, we can say that a decidable problem is one where we can always obtain an answer for both the problem itself and for the complement of the problem. Because both $\mathit{A}_{\DFA}$ and $\mathit{A}_{\CFG}$ are decidable, their Turing machines always halt and either accept or reject any input word given to them. This means that the corresponding ``non-membership" problems $\overline{\mathit{A}_{\DFA}}$ and $\overline{\mathit{A}_{\CFG}}$ are also decidable, since we can use the same Turing machines to decide these problems by exchanging accept and reject outputs, and so studying these complementary decision problems isn't too interesting. The Turing machine ``non-membership" problem $\overline{\mathit{A}_{\TM}}$, on the other hand, is different, since $\mathit{A}_{\TM}$ is only semidecidable. A semidecidable problem is one where we can always obtain an answer for the problem itself, but not necessarily for the complement of the problem. Thus, it makes sense for us to study $\overline{\mathit{A}_{\TM}}$ independently of $\mathit{A}_{\TM}$.

So, what is the decidability status of $\overline{\mathit{A}_{\TM}}$? Instead of approaching a proof from first principles, as we did in showing that $\mathit{A}_{\TM}$ was undecidable, we will use a different (and easier!) approach with $\overline{\mathit{A}_{\TM}}$ that focuses on what we already know about $\mathit{A}_{\TM}$.

Since Theorem~\ref{thm:ATMsemidecidable} tells us that $\mathit{A}_{\TM}$ is semidecidable, we know that $\overline{\mathit{A}_{\TM}}$ is co-semidecidable by Definition~\ref{def:cosemidecidablelanguage}. However, if $\overline{\mathit{A}_{\TM}}$ were also semidecidable, then Definition~\ref{def:cosemidecidablelanguage} would likewise tell us that $\mathit{A}_{\TM}$ is co-semidecidable, and this would present quite a problem! Therefore, we can conclude the following.

\begin{theorem}\label{thm:coATMnotsemidecidable}
$\overline{\mathit{A}_{\TM}}$ is not semidecidable.

\begin{proof}
Assume by way of contradiction that $\overline{\mathit{A}_{\TM}}$ is semidecidable. Then $\mathit{A}_{\TM}$ would be co-semidecidable, and by Theorem~\ref{thm:decidablesemidecidable}, $\mathit{A}_{\TM}$ would be decidable, which contradicts Theorem~\ref{thm:ATMundecidable}.
\end{proof}
\end{theorem}

This theorem reveals a remarkable fact: $\overline{\mathit{A}_{\TM}}$ is not only undecidable as we might expect, but it's also non-semidecidable, and so it lies in an entirely different part of our language hierarchy!

As we observed at the beginning of the section, uncountably many non-semidecidable languages exist, and so we shouldn't be too surprised by the fact that we found an example of such a language. In the next chapter, we'll see a number of other examples of languages that are neither decidable nor semidecidable; in some cases, these languages are co-semidecidable, while in other cases, not even co-semidecidability holds.

\unnumberedsection{Chapter Notes}

\begin{enumerate}
\item[\ref{sec:decidableREG}.] The decidability of the membership problem for the class of regular languages is practically folklore, in that the result follows immediately from simply giving an input word to a finite automaton. The decidability of the (non-)emptiness problem was established by \citet{Moore1956GedankenExperiments}, who used an argument different from ours that essentially resembles the pumping lemma for regular languages. As we observed, the decidability of the universality, equivalence, and inclusion problems all follow from the positive answer for the emptiness problem.

\item[\ref{sec:decidableCFL}.] \citet*{BarHillel1961FormalPropertiesPhraseStructureGrammars} showed that both the membership problem (or ``the property of sentencehood" in their terminology) and the emptiness problem are decidable for the class of context-free languages.

\item[\ref{sec:undecidableTM}.] It can be said that the undecidability of $\mathit{A}_{\TM}$ was established by \citet{Turing1936OnComputableNumbers}, although he didn't prove this exact result. Turing instead proved the undecidability of the \emph{printing problem}, which takes a machine $\mathcal{M}$ and asks ``whether $\mathcal{M}$ ever prints a given symbol (0 say)". We can, however, easily draw a connection between the printing problem and the membership problem by, for instance, modifying a given machine to write some special symbol ``\texttt{a}" to its tape before reaching $q_{\text{accept}}$.

Although Cantor's diagonal argument appeared in~\citeyear{Cantor1891Mannigfaltigkeitslehre}, it was not Cantor's first proof pertaining to uncountability. Two decades earlier, Cantor published a paper~\citeyearpar{Cantor1874Zahlen} demonstrating, among other things, the uncountability of the set of real numbers using an approach inspired by Liouville's argument for the existence of transcendental numbers. For further details about Cantor's proofs, see the article by~\citet{Gray1994CantorTranscendental}.

\item[\ref{sec:nonsemidecidableTM}.] The argument we employed in the proof of Theorem~\ref{thm:nonsemidecidablelanguages} is essentially identical to that used by \citet{Cantor1891Mannigfaltigkeitslehre}, but applied to languages instead of general sets.

Theorem~\ref{thm:decidablesemidecidable}, on the relationship between decidability, semidecidability, and co-semidecidability, is due to \citet{Post1944RecursivelyEnumerableSets}.
\end{enumerate}